% ===================== METODOLOGIA =====================
\section{Metodologia}
\label{sec:metodologia}

\subsection{Plan de Trabajo}

\begin{table}[htbp]
    \centering
    \caption{Plan de trabajo del an\'{a}lisis de flexibilidad.}
    \label{tab:plan_trabajo}
    \setcounter{itemcount}{0}
    \begin{tabularx}{\textwidth}{@{}NlX@{}}
        \toprule
        \multicolumn{1}{c}{\textbf{Etapa}} & \textbf{Actividad} & \textbf{Descripcion} \\
        \midrule
        & Modelado isom\'{e}trico & El equipo de dise\~{n}o genera el isom\'{e}trico PCF119 en AutoCAD Plant 3D 2027 y exporta el archivo de intercambio PCF. \\
        & Correcci\'{o}n del PCF & El archivo PCF se corrige con el script \texttt{fix\_pcf.py} para compatibilidad con CAESAR II 2019: eliminaci\'{o}n de bloques no soportados, fusi\'{o}n de \textit{stub-ends} con sus bridas y retiro de conexiones de instrumentaci\'{o}n. \\
        & Modelo de esfuerzos & El PCF corregido se importa en CAESAR II 2019; se definen las condiciones de operaci\'{o}n (secci\'{o}n~\ref{sec:bases}) y los casos de carga (Tabla~\ref{tab:casos}). \\
        & Verificaci\'{o}n y reporte & Se ejecuta el an\'{a}lisis est\'{a}tico y se verifica el resumen global de cumplimiento de c\'{o}digo; los reportes de desplazamientos, esfuerzos y restricciones se extraen del archivo de salida para este informe. \\
        \bottomrule
    \end{tabularx}
\end{table}

% ─────────────────────────────────────────────────────────────────────────────
\subsection{Herramientas de Analisis}

El modelo de esfuerzos se ejecut\'{o} en CAESAR II 2019 (Ver.11.00.00.4800,
Hexagon / Intergraph), software de an\'{a}lisis de flexibilidad basado en el
m\'{e}todo de elementos finitos con elementos viga tridimensionales de seis
grados de libertad por nodo (secci\'{o}n~\ref{ssec:marco})~\cite{caesar2_ug_2019}.
El isom\'{e}trico se produjo en AutoCAD Plant 3D 2027 (Autodesk) y la
transferencia de geometr\'{i}a se realiz\'{o} mediante el formato PCF,
previamente corregido con el script propio \texttt{fix\_pcf.py} (Python 3.14).
Las tablas y curvas de este informe se generaron program\'{a}ticamente desde el
reporte de salida de CAESAR II, sin transcripci\'{o}n manual de valores.

% ─────────────────────────────────────────────────────────────────────────────
\subsection{Matriz de Escenarios / Casos de Estudio}

Los nueve casos de carga evaluados se listan en la Tabla~\ref{tab:casos}
(secci\'{o}n~\ref{sec:bases}). El caso de expansi\'{o}n no
representa una condici\'{o}n f\'{i}sica de carga, sino el rango de
desplazamientos t\'{e}rmicos entre el estado de operaci\'{o}n y el estado
sostenido~\cite{whatispiping_pipe_stress}.
